\section{Dividing FPSs}

\subsection{Inverses in commutative rings}

\begin{definition}
\label{def.commring.inverse}
\lean{AlgebraicCombinatorics.inverse_unique}
\leantarget
\leanok
Let $L$ be a commutative ring. Let $a\in L$. Then:

\textbf{(a)} An \emph{inverse} (or \emph{multiplicative inverse}) of $a$ means
an element $b\in L$ such that $ab=ba=1$.

\textbf{(b)} We say that $a$ is \emph{invertible} in $L$ (or a \emph{unit} of
$L$) if $a$ has an inverse.
\end{definition}

\begin{theorem}
\label{thm.commring.inverse-uni}
\lean{AlgebraicCombinatorics.inverse_unique}
\leantarget
\leanok
Let $L$ be a commutative ring. Let $a\in L$.
Then, there is \textbf{at most one} inverse of $a$.
\end{theorem}

\begin{proof}
\leanok
Let $b$ and $c$ be two inverses of $a$.
Since $b$ is an inverse of $a$, we have $ab=ba=1$.
Since $c$ is an inverse of $a$, we have $ac=ca=1$.
Now, $b\left(ac\right) = b\cdot 1 = b$ and $\left(ba\right)c = 1\cdot c = c$.
However, $b\left(ac\right) = \left(ba\right)c$
by associativity. Hence, $b = b\left(ac\right) = \left(ba\right)c = c$.
\end{proof}

\begin{definition}
\label{def.commring.fracs}
\lean{AlgebraicCombinatorics.divByUnit}
\leantarget
\leanok
Let $L$ be a commutative ring. Let $a\in L$. Assume
that $a$ is invertible. Then:

\textbf{(a)} The inverse of $a$ is called $a^{-1}$.

\textbf{(b)} For any $b\in L$, the product $b\cdot a^{-1}$ is called
$\dfrac{b}{a}$ (or $b/a$).

\textbf{(c)} For any negative integer $n$, we define $a^{n}$ to be $\left(
a^{-1}\right)^{-n}$. Thus, the $n$-th power $a^{n}$ is defined for each
$n\in\mathbb{Z}$.
\end{definition}

\begin{proposition}
\label{prop.commring.fracs.1}
\lean{AlgebraicCombinatorics.fracs1_inv_eq_one_mul_inv, AlgebraicCombinatorics.fracs1_isUnit_mul, AlgebraicCombinatorics.fracs1_mul_inv_rev, AlgebraicCombinatorics.fracs1_mul_inv_comm, AlgebraicCombinatorics.fracs1_zpow_neg_eq_inv_zpow, AlgebraicCombinatorics.fracs1_zpow_add, AlgebraicCombinatorics.fracs1_mul_zpow, AlgebraicCombinatorics.fracs1_zpow_mul, AlgebraicCombinatorics.fracs1_div_add_div, AlgebraicCombinatorics.fracs1_div_mul_div, AlgebraicCombinatorics.fracs1_div_eq_iff}
\leantarget
\leanok
Let $L$ be a commutative ring. Then:

\textbf{(a)} Any invertible element $a\in L$ satisfies $a^{-1}=1/a$.

\textbf{(b)} For any invertible elements $a,b\in L$, the element $ab$ is
invertible as well, and satisfies $\left(ab\right)^{-1}=b^{-1}%
a^{-1}=a^{-1}b^{-1}$.

\textbf{(c)} If $a\in L$ is invertible, and if $n\in\mathbb{Z}$ is arbitrary,
then $a^{-n}=\left(a^{-1}\right)^{n}=\left(a^{n}\right)^{-1}$.

\textbf{(d)} Laws of exponents hold for negative exponents as well: If $a$ and
$b$ are invertible elements of $L$, then%
\begin{align*}
a^{n+m}  &  =a^{n}a^{m}\ \ \ \ \ \ \ \ \ \ \text{for all }n,m\in\mathbb{Z};\\
\left(  ab\right)  ^{n}  &  =a^{n}b^{n}\ \ \ \ \ \ \ \ \ \ \text{for all }%
n\in\mathbb{Z};\\
\left(  a^{n}\right)  ^{m}  &  =a^{nm}\ \ \ \ \ \ \ \ \ \ \text{for all
}n,m\in\mathbb{Z}.
\end{align*}

\textbf{(e)} Laws of fractions hold: If $a$ and $c$ are two invertible
elements of $L$, and if $b$ and $d$ are any two elements of $L$, then
\[
\dfrac{b}{a}+\dfrac{d}{c}=\dfrac{bc+ad}{ac}\ \ \ \ \ \ \ \ \ \ \text{and}%
\ \ \ \ \ \ \ \ \ \ \dfrac{b}{a}\cdot\dfrac{d}{c}=\dfrac{bd}{ac}.
\]

\textbf{(f)} Division undoes multiplication: If $a,b,c$ are three elements of
$L$ with $a$ being invertible, then the equality $\dfrac{c}{a}=b$ is
equivalent to $c=ab$.
\end{proposition}

\begin{proof}
Exercise (see, e.g., \cite[solution to Exercise 4.1.1]{19s} for a proof of
parts \textbf{(c)} and \textbf{(d)} in the special case where $L=\mathbb{C}$;
essentially the same argument works in the general case).
\end{proof}

\subsection{Inverses in $K[[x]]$}

\begin{proposition}
\label{prop.fps.invertible}
\lean{AlgebraicCombinatorics.fps_invertible_iff_constantCoeff}
\leantarget
\leanok
Let $a\in K[[x]]$.
Then, the FPS $a$ is invertible in $K[[x]]$ if
and only if its constant term $[x^{0}]a$ is invertible in $K$.
\end{proposition}

\begin{proof}
\leanok
$\Longrightarrow$: If $ab=1$, then $[x^0]a\cdot[x^0]b = [x^0](ab)=1$
(by Lemma~\ref{lem.fps.coeff.zero.mul}).

$\Longleftarrow$: If $[x^0]a = a_0$ is invertible, define a sequence
$(b_0, b_1, b_2, \ldots)$ recursively by
$b_0 = a_0^{-1}$ (Lemma~\ref{lem.fps.inv-coeff-zero}) and
$b_n = -a_0^{-1}\sum_{k=1}^n a_k b_{n-k}$ for $n \geq 1$
(Lemma~\ref{lem.fps.inv-coeff-succ}).
Then $ab = 1$.
\end{proof}

\begin{corollary}
\label{cor.fps.invertible.field}
\lean{AlgebraicCombinatorics.fps_invertible_iff_constantCoeff_ne_zero}
\leantarget
\leanok
Assume that $K$ is a field. Let $a\in K[[x]]$.
Then, the FPS $a$ is invertible in $K[[x]]$ if and only if $[x^{0}]a\neq0$.
\end{corollary}

\begin{proof}
\leanok
An element of $K$ is invertible in $K$ if and only if it is nonzero (since $K$
is a field). Hence, Corollary~\ref{cor.fps.invertible.field} follows from
Proposition~\ref{prop.fps.invertible}.
\end{proof}

\subsection{Coefficient formulas for inverses}

\begin{lemma}
\label{lem.fps.inv-coeff-zero}
\lean{AlgebraicCombinatorics.fps_inv_coeff_zero}
\leanhelper
\leanok
Assume that $K$ is a field. Let $f\in K[[x]]$. Then the constant term of $f^{-1}$
equals the inverse of the constant term of $f$:
\[
\left[x^{0}\right]f^{-1} = \left(\left[x^{0}\right]f\right)^{-1}.
\]
\end{lemma}

\begin{proof}
Direct consequence of the definition of the inverse of an FPS.
\end{proof}

\begin{lemma}
\label{lem.fps.inv-coeff-succ}
\lean{AlgebraicCombinatorics.fps_inv_coeff_succ}
\leanhelper
\leanok
Assume that $K$ is a field. Let $f\in K[[x]]$. For each $n\in\mathbb{N}$,
\[
\left[x^{n+1}\right]f^{-1}
= -\left(\left[x^{0}\right]f\right)^{-1}\sum_{k=0}^{n}\left[x^{k+1}\right]f\cdot
\left[x^{n-k}\right]f^{-1}.
\]
This recurrence allows computing the coefficients of $f^{-1}$ one by one.
\end{lemma}

\begin{proof}
Follows from the recurrence relation for the coefficients of $f^{-1}$ obtained
by expanding the equation $f\cdot f^{-1}=1$ and solving for $[x^{n+1}]f^{-1}$.
\end{proof}

\begin{lemma}
\label{lem.fps.isUnit-inv-eq-inv}
\lean{AlgebraicCombinatorics.fps_isUnit_inv_eq_inv}
\leanhelper
\leanok
Assume that $K$ is a field. Let $f\in K[[x]]$ be an invertible FPS.
Then the inverse of $f$ obtained from the invertibility proof equals
$f^{-1}$ (the inverse as defined in the power series ring).
\end{lemma}

\begin{proof}
\leanok
Both satisfy $f\cdot g = 1$, and inverses are unique
(Theorem~\ref{thm.commring.inverse-uni}).
\end{proof}

\begin{lemma}
\label{lem.fps.inv-coeff-zero-isUnit}
\lean{AlgebraicCombinatorics.fps_inv_coeff_zero_isUnit}
\leanhelper
\leanok
Assume that $K$ is a field. Let $f\in K[[x]]$ be an invertible FPS. Then
\[
\left[x^{0}\right]f^{-1} = \left(\left[x^{0}\right]f\right)^{-1}.
\]
\end{lemma}

\begin{proof}
\leanok
By Lemma~\ref{lem.fps.isUnit-inv-eq-inv}, the inverse obtained from
the invertibility proof equals $f^{-1}$,
so the result follows from Lemma~\ref{lem.fps.inv-coeff-zero}.
\end{proof}

\begin{lemma}
\label{lem.fps.inv-coeff-succ-isUnit}
\lean{AlgebraicCombinatorics.fps_inv_coeff_succ_isUnit}
\leanhelper
\leanok
Assume that $K$ is a field. Let $f\in K[[x]]$ be an invertible FPS. For each
$n\in\mathbb{N}$,
\[
\left[x^{n+1}\right]f^{-1}
= -\left(\left[x^{0}\right]f\right)^{-1}\sum_{k=0}^{n}\left[x^{k+1}\right]f\cdot
\left[x^{n-k}\right]f^{-1}.
\]
\end{lemma}

\begin{proof}
By Lemma~\ref{lem.fps.isUnit-inv-eq-inv}, the inverse obtained from
the invertibility proof equals $f^{-1}$,
so the result follows from Lemma~\ref{lem.fps.inv-coeff-succ}.
\end{proof}

\subsection{Newton's binomial formula}

\begin{proposition}
\label{prop.fps.invertible.1+x}
\lean{AlgebraicCombinatorics.fps_onePlusX_isUnit, AlgebraicCombinatorics.fps_onePlusX_inv}
\leantarget
\leanok
The FPS $1+x\in K[[x]]$ is invertible, and its inverse is%
\[
\left(1+x\right)^{-1}=\sum_{n\in\mathbb{N}}\left(-1\right)^{n}x^{n}.
\]
\end{proposition}

\begin{proof}
\leanok
Direct computation: $(1+x)\cdot\sum_{n\geq 0}(-1)^n x^n = 1$ by telescoping
(all powers of $x$ other than $1$ cancel out).
\end{proof}

\begin{theorem}
\label{thm.fps.newton-binom}
\lean{AlgebraicCombinatorics.fps_newtonBinomial_nat, AlgebraicCombinatorics.fps_newtonBinomial_neg}
\leantarget
\leanok
For each $n\in\mathbb{Z}$, we have%
\[
\left(1+x\right)^{n}=\sum_{k\in\mathbb{N}}\dbinom{n}{k}x^{k}.
\]
\end{theorem}

\begin{proof}
\leanok
For $n\geq 0$, this is the standard binomial theorem (the sum
$\sum_{k\in\mathbb{N}}\dbinom{n}{k}x^{k}$ reduces to $\sum_{k=0}^{n}\dbinom{n}{k}x^{k}$
since $\dbinom{n}{k}=0$ for $k > n$ by Proposition~\ref{prop.binom.0}).

For $n<0$, we have $-n\in\mathbb{N}$, so Corollary~\ref{cor.fps.anti-newton-binom-2}
(applied to $-n$ instead of $n$) yields
$\left(1+x\right)^{-(-n)} = \sum_{k\in\mathbb{N}}\dbinom{-(-n)}{k}x^k$,
which rewrites as $\left(1+x\right)^n = \sum_{k\in\mathbb{N}}\dbinom{n}{k}x^k$.
\end{proof}

\begin{theorem}
\label{thm.binom.upneg-n}
\lean{AlgebraicCombinatorics.binomUpperNegation, AlgebraicCombinatorics.binomUpperNegation_int}
\leantarget
\leanok
Let $n\in\mathbb{C}$ and $k\in\mathbb{Z}$. Then,%
\[
\dbinom{-n}{k}=\left(-1\right)^{k}\dbinom{k+n-1}{k}.
\]
\end{theorem}

\begin{proof}
\leanok
If $k<0$, then both $\dbinom{-n}{k}$ and $\dbinom{k+n-1}{k}$ are $0$.
For $k\geq 0$, expand using the definition and
observe the numerators differ by a factor of $(-1)^k$.
\end{proof}

\begin{proposition}
\label{prop.fps.anti-newton-binom}
\lean{AlgebraicCombinatorics.fps_onePlusX_pow_neg}
\leantarget
\leanok
For each $n\in\mathbb{N}$, we have%
\[
\left(1+x\right)^{-n}=\sum_{k\in\mathbb{N}}\left(-1\right)^{k}%
\dbinom{n+k-1}{k}x^{k}.
\]
\end{proposition}

\begin{proof}
By induction on $n$.

\textit{Induction base:} $\left(1+x\right)^{-0} = 1$, and
$\sum_{k\in\mathbb{N}} (-1)^k \dbinom{0+k-1}{k} x^k = 1$
since $\dbinom{k-1}{k} = 0$ for all $k > 0$ (by Proposition~\ref{prop.binom.0}).

\textit{Induction step:} Assume the formula holds for $n = j$. We must prove it for $n = j+1$.
We have $\left(1+x\right)^{-(j+1)} = \left(1+x\right)^{-j} \cdot \left(1+x\right)^{-1}$.
Using the induction hypothesis and multiplying by $1+x$, it suffices to show
\[
\left(1+x\right)^{-j} = \left(\sum_{k\in\mathbb{N}} (-1)^k \dbinom{j+k}{k} x^k\right) \cdot \left(1+x\right).
\]
Expanding the right hand side, using Pascal's identity
$\dbinom{j+k}{k} = \dbinom{j+k-1}{k-1} + \dbinom{j+k-1}{k}$
(Proposition~\ref{prop.binom.rec}), and collecting terms (the
$\dbinom{j+k-1}{k-1}$ terms telescope), we obtain
$\sum_{k\in\mathbb{N}} (-1)^k \dbinom{j+k-1}{k} x^k = \left(1+x\right)^{-j}$
by the induction hypothesis.
\end{proof}

\begin{corollary}
\label{cor.fps.anti-newton-binom-2}
\lean{AlgebraicCombinatorics.fps_onePlusX_pow_neg'}
\leantarget
\leanok
For each $n\in\mathbb{N}$, we have%
\[
\left(1+x\right)^{-n}=\sum_{k\in\mathbb{N}}\dbinom{-n}{k}x^{k}.
\]
\end{corollary}

\begin{proof}
\leanok
Proposition~\ref{prop.fps.anti-newton-binom} yields
$\left(1+x\right)^{-n} = \sum_{k\in\mathbb{N}} (-1)^k \dbinom{n+k-1}{k} x^k$.
By Theorem~\ref{thm.binom.upneg-n}, we have
$(-1)^k \dbinom{n+k-1}{k} = (-1)^k \dbinom{k+n-1}{k} = \dbinom{-n}{k}$.
\end{proof}

\subsection{Dividing by $x$}

\begin{definition}
\label{def.fps.div-by-x}
\lean{AlgebraicCombinatorics.PowerSeries.divByX}
\leantarget
\leanok
Let $a=\left(a_{0},a_{1},a_{2},\ldots\right)$ be
an FPS whose constant term $a_{0}$ is $0$. Then, $\dfrac{a}{x}$ is defined to
be the FPS $\left(a_{1},a_{2},a_{3},\ldots\right)$.
\end{definition}

\begin{lemma}
\label{lem.fps.coeff-divByX}
\lean{AlgebraicCombinatorics.PowerSeries.coeff_divByX}
\leanhelper
\leanok
Let $a\in K[[x]]$ with $[x^0]a=0$. Then for each $n\in\mathbb{N}$,
\[
\left[x^{n}\right]\frac{a}{x} = \left[x^{n+1}\right]a.
\]
\end{lemma}

\begin{proof}
\leanok
By definition, $\frac{a}{x}=(a_1,a_2,a_3,\ldots)$, so its $n$-th coefficient
is $a_{n+1}=[x^{n+1}]a$.
\end{proof}

\begin{proposition}
\label{prop.fps.div-by-x-inverts}
\lean{AlgebraicCombinatorics.fps_eq_X_mul_iff}
\leantarget
\leanok
Let $a,b\in K[[x]]$ be two FPSs. Then, $a=xb$ if
and only if $\left[x^{0}\right]a=0$ and $b=\dfrac{a}{x}$.
\end{proposition}

\begin{proof}
\leanok
Follows directly from the definitions.
\end{proof}

\subsubsection{Helpers for division by $x$}

\begin{lemma}
\label{lem.fps.eq-X-mul-divByX}
\lean{AlgebraicCombinatorics.fps_eq_X_mul_divByX}
\leanhelper
\leanok
Let $a\in K[[x]]$ with $[x^0]a=0$. Then $a = x\cdot\frac{a}{x}$.
\end{lemma}

\begin{proof}
\leanok
Follows from Proposition~\ref{prop.fps.div-by-x-inverts}: setting
$b=\frac{a}{x}$, the right-to-left direction gives $a=xb$.
\end{proof}

\begin{lemma}
\label{lem.fps.divByX-X-mul}
\lean{AlgebraicCombinatorics.fps_divByX_X_mul}
\leanhelper
\leanok
For any FPS $b\in K[[x]]$, we have $\frac{xb}{x}=b$.
\end{lemma}

\begin{proof}
\leanok
For each $n\in\mathbb{N}$, $[x^n]\frac{xb}{x}=[x^{n+1}](xb)=[x^n]b$
(by Lemma~\ref{lem.fps.coeff-divByX} and the coefficient formula for $xb$).
\end{proof}

\subsection{A few lemmas}

\begin{lemma}
\label{lem.fps.g=xh}
\lean{AlgebraicCombinatorics.fps_exists_X_mul_of_constantCoeff_zero}
\leantarget
\leanok
Let $a\in K[[x]]$ be an FPS
with $\left[x^{0}\right]a=0$. Then, there exists an $h\in K[[x]]$ such that $a=xh$.
\end{lemma}

\begin{proof}
\leanok
The FPS $\dfrac{a}{x}$ is well-defined (since $a_0 = 0$).
Set $h = \dfrac{a}{x}$; then $a = x \cdot \dfrac{a}{x} = xh$
(by Lemma~\ref{lem.fps.eq-X-mul-divByX}).
\end{proof}

\begin{lemma}
\label{lem.fps.first-n-coeffs-of-xna}
\lean{AlgebraicCombinatorics.fps_coeff_X_pow_mul_eq_zero}
\leantarget
\leanok
Let $k\in\mathbb{N}$. Let $a\in K[[x]]$ be any FPS. Then, the first $k$ coefficients of
the FPS $x^{k}a$ are $0$.
\end{lemma}

\begin{proof}
\leanok
For $m<k$, we have
$[x^m](x^k a) = \sum_{i=0}^{m} [x^i](x^k)\cdot [x^{m-i}]a = 0$
since $[x^i](x^k) = 0$ for all $i\leq m < k$.
\end{proof}

\begin{lemma}
\label{lem.fps.muls-of-xn}
\lean{AlgebraicCombinatorics.fps_first_k_coeffs_zero_iff_X_pow_dvd}
\leantarget
\leanok
Let $k\in\mathbb{N}$. Let $f\in K[[x]]$ be any FPS. Then, the first $k$ coefficients of
$f$ are $0$ if and only if $f$ is a multiple of $x^{k}$.
\end{lemma}

\begin{proof}
\leanok
$\Longleftarrow$: By Lemma~\ref{lem.fps.first-n-coeffs-of-xna}.

$\Longrightarrow$: If $f_n = 0$ for each $n \in \{0, 1, \ldots, k-1\}$, then
$f = \sum_{n \geq k} f_n x^n = x^k \sum_{n \geq k} f_n x^{n-k}$.
\end{proof}

\begin{lemma}
\label{lem.fps.prod.irlv.fg}
\lean{AlgebraicCombinatorics.fps_coeff_mul_eq_of_coeff_eq}
\leantarget
\leanok
Let $a,f,g\in K[[x]]$ be three FPSs. Let $n\in\mathbb{N}$. Assume that
$[x^{m}]f=[x^{m}]g$ for each $m\in\{0,1,\ldots,n\}$.
Then,
$[x^{m}](af)=[x^{m}](ag)$ for each $m\in\{0,1,\ldots,n\}$.
\end{lemma}

\begin{proof}
\leanok
Let $m \in \{0,1,\ldots,n\}$. Each $j \in \{0,1,\ldots,m\}$ satisfies $j \leq m \leq n$
and thus $[x^j]f = [x^j]g$. Hence,
$[x^m](af) = \sum_{j=0}^m [x^{m-j}]a \cdot [x^j]f
= \sum_{j=0}^m [x^{m-j}]a \cdot [x^j]g = [x^m](ag)$.
\end{proof}

\begin{lemma}
\label{lem.fps.prod.irlv.mul}
\lean{AlgebraicCombinatorics.fps_coeff_zero_of_multiple}
\leantarget
\leanok
Let $u,v\in K[[x]]$ be two FPSs such that $v$ is a multiple of $u$. Let $n\in\mathbb{N}$. Assume
$[x^{m}]u=0$ for each $m\in\{0,1,\ldots,n\}$.
Then $[x^{m}]v=0$ for each $m\in\{0,1,\ldots,n\}$.
\end{lemma}

\begin{proof}
\leanok
Write $v=ua$. We have $[x^m]u = 0 = [x^m]0$ for each $m \in \{0,1,\ldots,n\}$.
Lemma~\ref{lem.fps.prod.irlv.fg} (applied to $f=u$ and $g=0$) yields
$[x^m](au) = [x^m](a \cdot 0) = 0$ for each $m \in \{0,1,\ldots,n\}$.
Since $v = ua = au$, we get $[x^m]v = 0$.
\end{proof}

\begin{lemma}
\label{lem.fps.prod.irlv.cong-mul}
\lean{AlgebraicCombinatorics.fps_coeff_mul_eq_of_both_coeff_eq}
\leantarget
\leanok
Let $a,b,c,d\in K[[x]]$ be four FPSs. Let $n\in\mathbb{N}$. Assume
$[x^{m}]a=[x^{m}]b$ and $[x^{m}]c=[x^{m}]d$ for each $m\in\{0,1,\ldots,n\}$.
Then $[x^{m}](ac)=[x^{m}](bd)$ for each $m\in\{0,1,\ldots,n\}$.
\end{lemma}

\begin{proof}
\leanok
The FPS $ac - bc = (a-b)c$ is a multiple of $a-b$, and $[x^m](a-b) = 0$
for each $m \in \{0,1,\ldots,n\}$. By Lemma~\ref{lem.fps.prod.irlv.mul},
$[x^m](ac - bc) = 0$, so $[x^m](ac) = [x^m](bc)$.

Similarly, $bc - bd = b(c-d) = (c-d)b$ is a multiple of $c-d$, and
$[x^m](c-d) = 0$. By Lemma~\ref{lem.fps.prod.irlv.mul},
$[x^m](bc - bd) = 0$, so $[x^m](bc) = [x^m](bd)$.

Combining: $[x^m](ac) = [x^m](bc) = [x^m](bd)$.
\end{proof}
